\title{Direct Preference Optimization: Your Language Model is Secretly a Reward Model}

\begin{abstract}
Although large-scale unsupervised language models (LMs) acquire extensive world knowledge and certain reasoning capabilities, precisely steering their behavior remains challenging due to the entirely unsupervised nature of their training. Current approaches to achieve such control involve collecting human judgments on the relative quality of model outputs and fine-tuning the unsupervised LM to align with these preferences, frequently using reinforcement learning from human feedback (RLHF). Nevertheless, RLHF is a complicated and often unstable process that first fits a reward model to capture human preferences, and then fine-tunes the large unsupervised LM via reinforcement learning to maximize this learned reward while preventing excessive divergence from the original model. In this work, we propose a novel parameterization of the reward model within RLHF that allows for extracting the optimal policy in closed form, enabling us to address the standard RLHF task using a straightforward classification loss alone. The resulting method, which we term \textit{Direct Preference Optimization} (DPO), is stable, effective, and computationally efficient, removing the requirement for sampling from the LM during fine-tuning or extensive hyperparameter tuning. Our experiments demonstrate that DPO can fine-tune LMs to align with human preferences as well as or better than previous approaches. Importantly, fine-tuning with DPO surpasses PPO-based RLHF in controlling the sentiment of generated text, and matches or improves response quality in summarization and single-turn dialogue tasks while being considerably simpler to implement and train.
\end{abstract}

\section{Introduction}
Large-scale unsupervised language models (LMs) trained on extensive datasets develop remarkable capabilities~\citep{chowdhery2022palm, brown2020language, touvron2023llama,bubeck2023sparks}. However, these models learn from data created by humans who have diverse objectives, priorities, and skill levels. Some of these objectives and skills may be undesirable to replicate; for instance, although we want our AI coding assistant to \textit{recognize} common programming errors to fix them, when generating code, we prefer to steer the model toward the (possibly rare) high-quality coding skills embedded in its training data. Likewise, we might want our language model to be \textit{cognizant} of a widespread misconception held by half of the population, but we certainly do not want the model to assert this misconception as true in 50\% of its responses. Put differently, it is essential to select the model's \emph{preferred responses and behaviors} from its extensive \textit{knowledge and capabilities} to create AI systems that are safe, effective, and controllable \citep{ouyang2022training}. Although current approaches generally guide LMs to align with human preferences via reinforcement learning (RL), we will demonstrate that the RL-based objective commonly used can be exactly optimized through a straightforward binary cross-entropy objective, significantly simplifying the preference learning process.

Broadly speaking, existing techniques instill desired behaviors into a language model by leveraging curated datasets of human preferences that reflect behaviors deemed safe and beneficial by humans. This preference learning step follows an initial phase of large-scale unsupervised pre-training on a substantial text corpus. While the simplest method for preference learning is supervised fine-tuning on human demonstrations of high-quality responses, the most effective approaches employ reinforcement learning from human (or AI) feedback (RLHF/RLAIF; \citep{christiano2017deep,bai2022constitutional}). RLHF methods involve training a reward model on a dataset of human preferences and then using RL to optimize the language model policy to generate outputs that receive high reward while avoiding excessive deviation from the original model. Although RLHF yields models with impressive conversational and coding skills, its training pipeline is considerably more complex than supervised learning, requiring training multiple LMs and performing policy sampling during training, which results in substantial computational overhead.

In this work, we demonstrate a method to optimize a language model to conform to human preferences directly, without relying on explicit reward modeling or reinforcement learning. We introduce \textit{Direct Preference Optimization} (DPO), an approach that implicitly optimizes the same objective as current RLHF algorithms (reward maximization constrained by KL-divergence) while being easy to implement and train. Conceptually, the DPO update increases the relative log probability of preferred responses compared to dispreferred ones, but it includes a dynamic, example-specific importance weight that prevents model collapse observed with a naive probability ratio objective. Similar to existing methods, DPO depends on a theoretical preference model (such as the Bradley-Terry model; \cite{bradley1952rankanalysis}) to assess how well a reward function aligns with observed preference data. However, unlike prior approaches that use the preference model to define a preference loss for training a reward model followed by a policy optimizing that reward model, DPO employs a change of variables to express the preference loss as a direct function of the policy. Given a dataset of human preferences over model outputs, DPO thus optimizes a policy through a straightforward binary cross-entropy loss, yielding the optimal policy corresponding to an implicit reward function fitted to the preference data.

Our primary contribution is Direct Preference Optimization (DPO), a simple, reinforcement learning-free algorithm for training language models based on preferences. Experimental results demonstrate that DPO is at least as effective as existing techniques, including PPO-based RLHF, for preference learning in tasks such as sentiment adjustment, summarization, and dialogue, utilizing language models with up to 6B parameters.

\section{Related Work}

Self-supervised language models, as their scale increases, acquire the ability to perform certain tasks zero-shot \citep{radford2019language} or with only a few-shot prompts \citep{gpt3,megatron,chowdhery2022palm}. Nevertheless, their effectiveness on downstream applications and alignment with user intentions can be greatly enhanced by fine-tuning on instruction datasets paired with human-generated completions \citep{mishra-etal-2022-cross,sanh2022multitask,chung2022scaling,thoppilan2022lamda}. This process, known as `instruction-tuning,' allows large language models (LLMs) to generalize beyond the specific instruction-tuning data and typically improves their overall usability \citep{chung2022scaling}. Although instruction tuning has proven successful, collecting \textit{relative} human judgments of response quality is often simpler than obtaining expert demonstrations. Consequently, later research has fine-tuned LLMs using datasets based on human preferences, resulting in enhanced performance in tasks such as translation \citep{kreutzer-etal-2018-reliability}, summarization \citep{stiennon2022learning,ziegler2020finetuning}, storytelling \citep{ziegler2020finetuning}, and instruction adherence \citep{ouyang2022training,ramamurthy2023is}. These approaches first train a neural network reward function to align with the preference dataset under models like the Bradley-Terry model \citep{bradley1952rankanalysis}, then fine-tune the language model to maximize this reward using reinforcement learning techniques, commonly REINFORCE \citep{williams1992reinforce}, proximal policy optimization (PPO; \cite{schulman2017proximal}), or related variants \citep{ramamurthy2023is}. A related research direction utilizes LLMs fine-tuned for instruction following with human feedback to create additional synthetic preference data focused on specific attributes such as safety or harmlessness \citep{bai2022constitutional}, relying solely on weak supervision from humans via textual rubrics for the LLM’s annotations. These strategies represent a merging of two research domains: one focused on training language models with reinforcement learning across various objectives~\citep{Ranzato2015SequenceLT,paulus2018a,wu2018learning} and another on general methodologies for learning from human preferences \citep{christiano2017deep,kupcsik2018learning}. Despite the attractiveness of utilizing relative human preferences, fine-tuning large language models through reinforcement learning remains a significant practical hurdle; this paper proposes a theoretically grounded method to optimize relative preferences without relying on RL.

Beyond the language domain, learning policies from preferences has been explored in both bandit and reinforcement learning frameworks, with various methods proposed. Contextual bandit learning based on preferences or rankings of actions, instead of rewards, is referred to as a contextual dueling bandit (CDB; \cite{yue2012karmed,dudik2015contextual}). When absolute rewards are unavailable, theoretical analyses of CDBs replace the concept of an optimal policy with a \textit{von Neumann winner}, a policy whose expected probability of winning against \textit{any} other policy is at least 50\% \citep{dudik2015contextual}. However, in the CDB framework, preference annotations are collected online, whereas learning from human preferences generally involves training on a fixed batch of offline preference-annotated action pairs \citep{yan2022human}. Similarly, \textit{preference-based RL} (PbRL) involves learning from binary preferences generated by an \textit{unknown} `scoring' function instead of explicit rewards \citep{BusaFekete2014,ruiz2023dueling}. Multiple PbRL algorithms exist, including those capable of reusing off-policy preference data, but they typically require first explicitly estimating the underlying scoring function (i.e., the reward model) before optimizing it \citep{jain2013learning,BusaFekete2014,christiano2017deep,sadigh2017active,kupcsik2018learning}. In contrast, we propose a unified, single-stage policy learning method that directly optimizes a policy to align with preferences.

\section{Preliminaries}\label{section:prelims}

We summarize the RLHF pipeline as described by \citeauthor{ziegler2020finetuning} and later works \citep{stiennon2022learning, bai2022training, ouyang2022training}. This typically comprises three steps: 1) supervised fine-tuning (SFT); 2) preference collection and reward modeling; and 3) RL-based optimization.

\textbf{SFT}: RLHF commonly starts by fine-tuning a pre-trained language model (LM) on high-quality, task-specific datasets (e.g., dialogue, summarization) via supervised learning, resulting in a model denoted $\pi^\text{SFT}$.

\textbf{Reward Modeling Phase}: In the second step, the SFT model is prompted with inputs $x$ to generate pairs of responses $(y_1, y_2) \sim \pi^\text{SFT}(y \mid x)$. These pairs are shown to human annotators, who indicate their preference for one answer, represented as $y_w \succ y_l \mid x$, where $y_w$ is the preferred completion and $y_l$ is the less preferred one among $(y_1, y_2)$. These preferences are assumed to arise from some latent reward function $r^*(y, x)$, which is unobserved. Various models can be employed to capture preferences, with the Bradley-Terry (BT) model \cite{bradley1952rankanalysis} being a common choice (although more general Plackett-Luce ranking models \citep{plackett1975analysis, luce2012individual} can also be used if multiple ranked answers are available). The BT model specifies that the human preference distribution $p^*$ can be expressed as:

\begin{equation}\label{eq:bradley-terry}
    p^*(y_1\succ y_2 \mid x)=\frac{\exp\left(r^*(x, y_1)\right)}{\exp\left(r^*(x, y_1)\right) + \exp\left(r^*(x, y_2)\right)}.
\end{equation}

Given access to a fixed dataset of comparisons $\mathcal{D}=\bigl\{x^{(i)}, y_w^{(i)}, y_l^{(i)}\bigr\}_{i=1}^N$ drawn from $p^*$, one can parameterize a reward model $r_{\phi}(x, y)$ and learn its parameters by maximizing the likelihood. Casting the task as a binary classification, the negative log-likelihood loss is formulated as:

\begin{equation}\label{eq:reward_model}
    \mathcal{L}_R(r_{\phi}, \mathcal{D}) = -\mathbb{E}_{(x, y_w, y_l)\sim \mathcal{D}}\bigl[\log \sigma(r_{\phi}(x, y_w)- r_{\phi}(x, y_l))\bigr]
\end{equation}

where $\sigma$ denotes the logistic function. For language models, $r_{\phi}(x, y)$ is often initialized based on the SFT model $\pi^\text{SFT}(y \mid x)$, enhanced by adding a linear layer atop the final transformer layer that outputs a single scalar reward prediction \cite{ziegler2020finetuning}. To reduce the variance of the reward function, previous works normalize the rewards such that $\mathbb{E}_{x,y\sim \mathcal{D}}\left[r_\phi(x, y)\right] = 0$ for every $x$.

\textbf{RL Fine-Tuning Phase}: In the reinforcement learning stage, the learned reward function guides the language model’s improvement. Following prior studies~\citep{jaques2017sequence, jaques2020human}, the optimization problem is expressed as:

\begin{equation}\label{eq:RL}
\max_{\pi_{\theta}}  \mathbb{E}_{x\sim \mathcal{D}, y\sim \pi_{\theta}(y \mid x)}\bigl[r_{\phi}(x, y)\bigr] - \beta\mathbb{D}_{\textrm{KL}}\bigl[\pi_{\theta}(y\mid x)\mid \mid \pi_\text{ref}(y\mid x)\bigr],
\end{equation}

where $\beta$ controls the extent of deviation from the baseline reference policy $\pi_\text{ref}$, typically the initial SFT model $\pi^\text{SFT}$. Practically, the policy $\pi_\theta$ of the language model is initialized to $\pi^\text{SFT}$. This constraint is crucial as it restricts the model from straying too far from the distribution where the reward model is valid, while also preserving generation diversity and avoiding mode collapse to a few high-reward outputs. Due to the discrete nature of language generation, this objective is non-differentiable and is generally optimized using reinforcement learning. The common approach \citep{ziegler2020finetuning, stiennon2022learning, bai2022training, ouyang2022training} involves defining the reward as ${r(x, y) = r_{\phi}(x, y) -\beta (\log \pi_{\theta}(y\mid x) - \log \pi_\text{ref}(y\mid x))}$ and maximizing it with PPO \cite{schulman2017proximal}.

\section{Direct Preference Optimization}\label{sec:DPO}

Inspired by the difficulties in applying reinforcement learning algorithms to large-scale tasks such as fine-tuning language models, our objective is to develop a straightforward policy optimization method that directly utilizes preferences. Unlike prior RLHF approaches, which first learn a reward and then optimize it through RL, our method exploits a specific reward model parameterization that allows the optimal policy to be derived analytically in closed form, eliminating the need for an RL training loop. As we will elaborate shortly, the central idea is to use an analytical transformation from reward functions to optimal policies, enabling conversion of a loss defined over reward functions into one defined over policies. This change-of-variables technique bypasses the necessity to fit an explicit, standalone reward model, yet still facilitates optimization under established human preference models like the Bradley-Terry model. Essentially, the policy network simultaneously represents both the language model and the implicit reward.

\textbf{Deriving the DPO objective.} We begin with the same RL objective as in previous work, Eq.~\ref{eq:RL}, defined under a general reward function $r$. Building on prior studies~\citep{peters2007reinforcement, peng2019advantage, korbak2022reinforcement, go2023aligning}, it is straightforward to demonstrate that the optimal solution to the KL-constrained reward maximization problem in Eq.~\ref{eq:RL} is given by:

\begin{equation}\label{eq:op_policy}
    \pi_r(y\mid x) = \frac{1}{Z(x)}\pi_\text{ref}(y\mid x)\exp\left(\frac{1}{\beta}r(x, y)\right),
\end{equation}

where $Z(x) =\sum_{y}\pi_\text{ref}(y\mid x)\exp\left(\frac{1}{\beta}r(x, y)\right)$ denotes the partition function. For a detailed derivation, refer to Appendix \ref{app:derivation1}. Even if we replace $r$ with its MLE estimate $r_{\phi}$ of the true reward $r^*$, estimating $Z(x)$ remains computationally demanding \citep{korbak2022reinforcement, go2023aligning}, which complicates practical use of this formulation. Nevertheless, we can rearrange Eq.~\ref{eq:op_policy} to rewrite the reward function using the associated optimal policy $\pi_r$, the reference policy $\pi_\text{ref}$, and the unknown partition function $Z(\cdot)$. Specifically, by taking the logarithm of both sides of Eq.~\ref{eq:op_policy} and performing some algebraic manipulation, we obtain:

\begin{equation}\label{eq:main_eq}
    r(x,y) =\beta \log \frac{\pi_r(y\mid x)}{\pi_\text{ref}(y\mid x)} + \beta \log Z(x).
\end{equation}

This reparameterization can be applied to the ground-truth reward $r^*$ and its corresponding optimal model $\pi^*$. Importantly, the Bradley-Terry model depends solely on the reward difference between two completions, i.e., ${p^*(y_1 \succ y_2 \mid x) = \sigma(r^*(x, y_1) - r^*(x, y_2))}$. Substituting the expression for $r^*(x,y)$ from Eq.~\ref{eq:main_eq} into the preference model Eq.~\ref{eq:bradley-terry} causes the partition function terms to cancel out, allowing us to represent the human preference probability purely in terms of the optimal policy $\pi^*$ and the reference policy $\pi_\text{ref}$. Consequently, the optimal RLHF policy $\pi^*$ under the Bradley-Terry framework satisfies the preference model:

\begin{equation}\label{eq:objective}
    p^*(y_1\succ y_2 \mid x)=\frac{1}{1 + \exp\left(\beta \log \frac{\pi^*(y_2\mid x)}{\pi_\text{ref}(y_2\mid x)} - \beta \log \frac{\pi^*(y_1\mid x)}{\pi_\text{ref}(y_1\mid x)}\right)}
\end{equation}

The full derivation is provided in Appendix~\ref{app:derivation2}. While Eq.~\ref{eq:objective} is derived using the Bradley-Terry model, similar derivations can be made for more general Plackett-Luce models~\citep{plackett1975analysis, luce2012individual}, as detailed in Appendix~\ref{app:plackett_luce_models}.

Having expressed the probability of human preference data in terms of the optimal policy rather than the reward function, we can now define a maximum likelihood objective for a parameterized policy $\pi_\theta$. In analogy to the reward modeling approach (cf. Eq.~\ref{eq:reward_model}), the resulting policy objective is:

\begin{equation}\label{eq:optimum_model}
    \mathcal{L}_\text{DPO}(\pi_{\theta}; \pi_\text{ref}) = -\mathbb{E}_{(x, y_w, y_l)\sim \mathcal{D}}\left[\log \sigma \left(\beta \log \frac{\pi_{\theta}(y_w\mid x)}{\pi_\text{ref}(y_w\mid x)} - \beta

In this manner, we fit an implicit reward through an alternative parameterization, whose optimal policy corresponds directly to $\pi_\theta$. Additionally, because our method is analogous to fitting a reparameterized Bradley-Terry model, it benefits from certain theoretical guarantees, including consistency under appropriate assumptions about the distribution of preference data \cite{bong2022generalized}. We elaborate further on the theoretical aspects of DPO and its relationship to other works in Section~\ref{sec:theory}.

\textbf{What effect does the DPO update have?} To gain a mechanistic understanding of DPO, it is helpful to examine the gradient of the loss function $\mathcal{L}_\text{DPO}$. The gradient with respect to parameters $\theta$ is expressed as:

\begin{multline*}\label{eq:gradient}
    \nabla_\theta \mathcal{L}_\text{DPO}(\pi_\theta;\pi_\text{ref}) = \\ -\beta\mathbb{E}_{(x, y_w, y_l) \sim \mathcal{D}} \bigg[\underbrace{\sigma(\hat{r}_\theta(x, y_l) - \hat{r}_\theta (x, y_w))}_\text{greater weight when reward estimate is incorrect}\bigg[\underbrace{\nabla_\theta\log \pi(y_w \mid x)}_\text{boost likelihood of $y_w$} - \underbrace{\nabla_\theta\log\pi(y_l \mid x)}_\text{reduce likelihood of $y_l$}\bigg]\bigg],
\end{multline*}

where $\hat{r}_\theta(x, y) = \beta \log \frac{\pi_\theta(y \mid x)}{\pi_\text{ref}(y \mid x)}$ denotes the reward implicitly defined by the language model $\pi_\theta$ and the reference model $\pi_\text{ref}$ (further details in Section~\ref{sec:theory}). Intuitively, the gradient of $\mathcal{L}_\text{DPO}$ encourages increasing the probability of preferred completions $y_w$ and decreasing the probability of less preferred completions $y_l$. Crucially, examples are weighted by how much the implicit reward model $\hat{r}_\theta$ overestimates the dispreferred completions, scaled by $\beta$—that is, by how inaccurately the implicit reward model ranks the completions, reflecting the influence of the KL divergence constraint. Our experimental results highlight the significance of this weighting, as a naive variant of the method without this weighting factor can cause the language model to deteriorate (see Appendix Table~\ref{tab:unlikelihood_generations}).

\textbf{DPO outline.} The typical DPO procedure is as follows: 1) For each prompt $x$, sample completions $y_1, y_2 \sim \pi_\text{ref}(\cdot \mid x)$, label them using human preferences to form the offline preference dataset $\mathcal{D} = \{x^{(i)}, y_w^{(i)}, y_l^{(i)}\}_{i=1}^N$, and 2) train the language model $\pi_\theta$ to minimize the loss $\mathcal{L}_\text{DPO}$ given the fixed $\pi_\text{ref}$, dataset $\mathcal{D}$, and target parameter $\beta$. In practical applications, it is preferable to leverage publicly accessible preference datasets instead of generating samples and collecting human judgments anew. Because these datasets are generated using $\pi^\text{SFT}$, we set $\pi_\text{ref} = \pi^\text{SFT}$ when this model is available. In cases where $\pi^\text{SFT}$ is not accessible, we initialize $\pi_\text{ref}$ by maximizing the likelihood over preferred completions $(x, y_w)$, i.e., $\pi_\text{ref} = \argmax_{\pi} \mathbb{E}_{x,y_w \sim \mathcal{D}}[\log \pi(y_w \mid x)]$. This approach helps to reduce distributional discrepancies between the true, unknown reference distribution and the $\pi_\text{ref}$ employed in DPO. Additional specifics about implementation and hyperparameters are provided in Appendix~\ref{app:implementation}.

\section{Theoretical Analysis of DPO} This section offers a deeper interpretation of the DPO method, presents theoretical justification, and connects the benefits of DPO to challenges encountered in actor-critic algorithms applied to RLHF, such as PPO~\cite{schulman2017proximal}.

\label{sec:theory}

\subsection{Your Language Model Is Secretly a Reward Model} DPO avoids explicitly fitting a reward function and conducting reinforcement learning by utilizing a single maximum likelihood objective to learn the policy. Note that the optimization objective in Eq. \ref{eq:main_eq} corresponds to a Bradley-Terry model with a reward parameterization $r^*(x,y) = \beta \log \frac{\pi^*_\theta(y \mid x)}{\pi_\text{ref}(y \mid x)}$. We optimize the parametric model $\pi_{\theta}$ equivalently to the reward model optimization in Eq. \ref{eq:reward_model} via a change of variables. In this section, we develop the theoretical foundations for this reparameterization, demonstrate that it does not restrict the space of learnable reward functions, and show that it permits exact recovery of the optimal policy. We start by defining an equivalence relation among reward functions.

\begin{definition}
Two reward functions $r(x, y)$ and $r'(x, y)$ are considered equivalent if and only if there exists some function $f$ such that ${r(x, y)-r'(x, y) = f(x)}$.     
\end{definition}

It is straightforward to verify that this defines an equivalence relation, thereby dividing the set of reward functions into equivalence classes. We present the following two lemmas:

\begin{lemma}\label{lemma:same_prefrence} Within the Plackett-Luce model, and specifically the Bradley-Terry framework, any two reward functions belonging to the same equivalence class generate identical preference distributions.
\end{lemma}

\begin{lemma}\label{lemma:same_policy}
    Any two reward functions from the same equivalence class result in the same optimal policy when solving the constrained RL problem.
\end{lemma}

The proofs are simple and are provided in Appendix \ref{app:lemma1}. The first lemma captures a well-known under-identification issue inherent in the Plackett-Luce family of models \cite{plackett1975analysis}. Due to this ambiguity, additional identifiability constraints are typically required to ensure guarantees for the MLE estimators derived from Eq. \ref{eq:reward_model} \cite{bong2022generalized}. The second lemma implies that all reward functions within the same equivalence class correspond to the same optimal policy, so for our ultimate goal, it suffices to recover any one reward function from the optimal class. We establish the following theorem in Appendix~\ref{app:thm1}:

\begin{theorem}\label{thm:main}
    Under mild conditions, every reward equivalence class compatible with the Plackett-Luce (and specifically the Bradley-Terry) models can be expressed via the reparameterization ${r(x, y) = \beta \log \frac{\pi(y\mid x)}{\pi_\text{ref}(y\mid x)}}$ for some model $\pi(y\mid x)$ and a fixed reference model $\pi_\text{ref}(y \mid x)$.
\end{theorem}

\begin{sproof}
    Take any reward function $r(x, y)$ inducing an associated optimal policy $\pi_r(y \mid x)$ as defined in Eq. \ref{eq:op_policy}. We demonstrate that a reward function from $r$’s equivalence class can be represented in the form given above. Define the projection operator $f$ as  
\begin{equation}
    f(r; \pi_\text{ref}, \beta)(x, y) = r(x, y) - \beta\log\sum_{y}\pi_\text{ref}(y\mid x)\exp\left(\frac{1}{\beta}r(x, y)\right)
\end{equation}
This operator normalizes the reward function by subtracting the log of the partition function induced by $\pi_r$. Since this normalization term depends only on the prefix $x$, $f(r; \pi_\text{ref}, \beta)(x, y)$ remains within the equivalence class of $r(x, y)$. Substituting $r$ with the right-hand side of Eq.~\ref{eq:main_eq} (which is valid for any reward), it follows that $f(r; \pi_\text{ref}, \beta)(x, y) = \beta \log \frac{\pi_r(y\mid x)}{\pi_\text{ref}(y\mid x)}$. Hence, the projection $f$ yields a reward function in the equivalence class of $r$ that has the desired parameterization, showing that our proposed reparameterization does not reduce generality in modeling reward functions.
\end{sproof}

An alternative perspective on Theorem~\ref{thm:main} is that it precisely identifies which reward function within each equivalence class is chosen by the DPO reparameterization. Specifically, this is the reward function fulfilling:

\begin{equation}\label{eq:lag_p}
     \sum_{y}\underbrace{\pi_\text{ref}(y\mid x)\exp\left(\frac{1}{\beta}r(x, y)\right)}_{=\pi(y\mid x)\text{, using Thm.~\ref{thm:main} reparam.}} = 1,
\end{equation}

meaning that $\pi(y\mid x)$ constitutes a valid probability distribution (all probabilities are non-negative and sum to unity). However, by referring to Eq.~\ref{eq:op_policy}, we observe that Eq.~\ref{eq:lag_p} corresponds to the partition function of the optimal policy derived from the reward function $r(x, y)$. The fundamental insight underlying the DPO algorithm is that we can enforce specific constraints on the otherwise underdetermined Plackett-Luce family of preference models (notably the Bradley-Terry model), thereby maintaining the set of representable reward functions but making the optimal policy defined in Eq.~\ref{eq:op_policy} analytically manageable for any given prompt $x$.

\subsection{Instability of Actor-Critic Algorithms} Our framework also enables the diagnosis of stability issues inherent in conventional actor-critic methods commonly employed in RLHF, such as PPO. Following the RLHF procedure, we concentrate on the RL fine-tuning phase described in Section \ref{section:prelims}. We can establish connections to the control-as-inference framework \cite{levine2018reinforcement} as applied to the constrained RL problem formulated in \ref{eq:RL}. Assuming a parameterized policy model $\pi_{\theta}(y\mid x)$, the goal is to minimize the KL divergence $\mathbb{D}_{\text{KL}}[\pi_{\theta}(y|x) \mid \mid \pi^*(y\mid x)]$, where $\pi^*$ denotes the optimal policy given by Eq.~\ref{eq:optimum_model} based on the reward function $r_{\phi}(y, x)$. Through algebraic manipulation, this results in the following optimization objective:

\begin{equation}\label{eq:AC}
    \max_{\pi_{\theta}}\mathbb{E}_{\pi_{\theta}(y\mid x)}\bigg[\underbrace{r_{\phi}(x, y) -\beta\log\sum_{y}\pi_\text{ref}(y\mid x)\exp\left(\frac{1}{\beta}r_{\phi}(x, y)\right)}_{f(r_{\phi}, \pi_\text{ref}, \beta)} - \underbrace{\beta\log\frac{\pi_{\theta}(y\mid x)}{\pi_\text{ref}(y\mid x)}}_{\text{KL}}\bigg]
\end{equation}

This objective is identical to the one optimized in previous studies \citep{ziegler2020finetuning, stiennon2022learning, bai2022training, ouyang2022training}, which utilize the DPO-equivalent reward for the reward class $r_{\phi}$. In this context, the normalization term within $f(r_{\phi}, \pi_\text{ref}, \beta)$ can be viewed as the soft value function corresponding to the reference policy $\pi_\text{ref}$. Although this term does not influence the optimal solution, omitting it can lead to a policy gradient with high variance, thereby destabilizing learning. While it is possible to incorporate the normalization term via a learned value function, this approach can also present optimization challenges. Alternatively, previous works have normalized rewards using a human completion baseline, essentially implementing a single-sample Monte Carlo estimate of the normalization term. In contrast, the DPO reparameterization produces a reward function that eliminates the need for any baselines.

\section{Experiments}
In this section, we conduct an empirical evaluation of DPO's capacity to train policies directly from preference data. Initially, within a controlled text-generation environment, we investigate how efficiently DPO balances maximizing reward against minimizing KL-divergence relative to the reference policy, in comparison to widely used preference learning methods such as PPO. Subsequently, we assess DPO’s effectiveness on larger models and more challenging RLHF tasks, including summarization and dialogue. Our findings indicate that, with minimal hyperparameter tuning, DPO generally matches or surpasses the performance of strong baselines like RLHF with PPO, as well as outperforming the best among $N$ sampled trajectories under a learned reward function. Prior to discussing these results, we outline the experimental setup; further information is provided in Appendix~\ref{app:exp_details}.

\textbf{Tasks.} Our study investigates three distinct open-ended text generation tasks. For each experiment, algorithms learn a policy from a preference dataset $\mathcal{D}=\bigl\{x^{(i)}, y_w^{(i)}, y_l^{(i)}\bigr\}_{i=1}^N$. In \textbf{controlled sentiment generation}, $x$ represents a prefix of a movie review drawn from the IMDb dataset \cite{maas-EtAl:2011:ACL-HLT2011}, and the policy is required to generate $y$ that expresses a positive sentiment. To enable a controlled evaluation, we \textit{create} preference pairs over generated outputs in this experiment by employing a pre-trained sentiment classifier, such that $p(\text{positive}\mid x,y_w) > p(\text{positive}\mid x,y_l)$. For SFT, we fine-tune GPT-2-large on the training split of the IMDb reviews until convergence (additional details are provided in App~\ref{app:sentiment_details}). In \textbf{summarization}, $x$ is a Reddit forum post; the policy must generate a summary $y$ capturing the post's key points. Building on previous work, we utilize the Reddit TL;DR summarization dataset \citep{volske-etal-2017-tl} along with human preference annotations collected by \citeauthor{stiennon2022learning}. We employ an SFT model fine-tuned on human-written forum post summaries\footnote{\url{https://huggingface.co/CarperAI/openai_summarize_tldr_sft}} within the TRLX \citep{leandro_von_werra_2023_7790115} framework for RLHF. The human preference dataset was obtained by \citeauthor{stiennon2022learning} from samples generated by a different yet similarly-trained SFT model. Lastly, in \textbf{single-turn dialogue}, $x$ is a user query, which can range from astrophysics questions to relationship advice requests. The policy must generate a response $y$ that is both engaging and helpful. We use the Anthropic Helpful and Harmless dialogue dataset \citep{bai2022training}, which includes 170k dialogues between humans and an automated assistant. Each dialogue ends with two responses generated by a large (though undisclosed) language model, accompanied by a preference label indicating the human-preferred response. In this task, no pre-trained SFT model is available; consequently, we fine-tune a readily available language model exclusively on the preferred completions to obtain the SFT model.

\textbf{Evaluation.} Our experiments employ two distinct evaluation methodologies. To assess how well each algorithm optimizes the constrained reward maximization objective in the controlled sentiment generation context, we evaluate each algorithm based on its achieved reward and KL-divergence from the reference policy, forming a frontier; this is feasible since the ground-truth reward function (a sentiment classifier) is available. In contrast, in real-world scenarios where the true reward function is unknown, we assess algorithms by their \textit{win rate} compared to a baseline policy, using GPT-4 as a stand-in for human judgment regarding summary quality and response helpfulness in the summarization and single-turn dialogue tasks, respectively. For summarization, the baseline consists of reference summaries from the test set; for dialogue, the baseline is the preferred response recorded in the test dataset. Although prior work indicates that language models can serve as better automated evaluators than standard metrics \citep{Chen2023ExploringTU}, we conduct a human study to validate the use of GPT-4 for evaluation as detailed in Sec.~\ref{sec:human-judgments}. Our findings reveal that GPT-4’s assessments strongly correlate with human evaluations, with agreement between humans and GPT-4 typically on par with or exceeding the agreement among different human annotators.

\textbf{Methods.} Beyond DPO, we assess various established techniques for training language models to align with human preferences. At the simplest level, we examine zero-shot prompting with \textbf{GPT-J} \citep{gpt-j} in summarization, and 2-shot prompting with \textbf{Pythia-2.8B} \citep{biderman2023pythia} in the dialogue task. We also evaluate the \textbf{SFT} model along with \textbf{Preferred-FT}, a model fine-tuned via supervised learning on the selected completion $y_w$ derived either from the SFT model (in controlled sentiment and summarization) or from a generic language model (in single-turn dialogue). Another pseudo-supervised approach is \textbf{Unlikelihood}~\citep{welleck2019neural}, which optimizes the policy by increasing the likelihood of $y_w$ while \textit{decreasing} the likelihood of $y_l$; this method optionally includes a coefficient $\alpha\in[0,1]$ scaling the ‘unlikelihood’ term. Additionally, we consider \textbf{PPO} \citep{schulman2017proximal}, which utilizes a reward function learned from preference data, and \textbf{PPO-GT}, an oracle version that learns directly from the ground truth reward function available in the controlled sentiment setting. For the sentiment experiments, we implement two variants of PPO-GT: an off-the-shelf version \cite{leandro_von_werra_2023_7790115} and a modified version that normalizes rewards and fine-tunes hyperparameters for enhanced performance (these modifications are also applied to the standard PPO with learned rewards). Lastly, we include the \textbf{Best of $N$} baseline, which samples $N$ responses from the SFT model (or Preferred-FT in dialogue) and selects the highest-scoring response based on a reward model trained from the preference dataset. While this method often yields strong performance by separating the quality of the reward model from PPO optimization, it is computationally expensive for even moderate values of $N$, as it requires generating $N$ completions per query during evaluation.

\subsection{To what extent can DPO optimize the RLHF objective?}

The KL-constrained reward maximization goal employed in standard RLHF methods strikes a balance between maximizing reward and limiting the policy’s deviation from the reference policy. Consequently, when evaluating different algorithms, it is essential to consider both the reward attained and the KL divergence; obtaining a marginally higher reward at the cost of significantly increased KL divergence is not necessarily preferable. Figure~\ref{fig:frontier-tldr-main} illustrates the reward-KL trade-off curve for various algorithms within the sentiment task. We perform multiple training runs per algorithm, each run varying the hyperparameter controlling policy conservativeness (target KL values $\in \{3,6,9,12\}$ for PPO, $\beta \in \{0.05,0.1,1,5\}$, $\alpha \in \{0.05,0.1,0.5,1\}$ for unlikelihood, and different random seeds for preferred-FT). This grid search comprises a total of 22 runs. At intervals of every 100 training steps until convergence, each policy is evaluated on a set of test prompts by calculating the average reward according to the true reward function as well as the average sequence-level KL\footnote{Specifically, the sum of the KL divergences at each timestep.} with respect to the reference policy, $\text{KL}\left(\pi \mid\mid \pi_\text{ref}\right)$. Our findings show that DPO yields the most efficient frontier by a large margin, delivering the highest reward while maintaining low KL divergence. This outcome is especially significant for several reasons. First, although DPO and PPO optimize the same objective, DPO is considerably more efficient; the reward-KL trade-off for DPO consistently surpasses that of PPO. Second, DPO outperforms PPO even when PPO is provided access to ground truth rewards (PPO-GT).

\subsection{Is DPO scalable to real preference datasets?}
\label{sec:dpo-real-datasets}
Next, we assess the fine-tuning effectiveness of DPO on tasks involving summarization and single-turn dialogue. In summarization, automatic metrics like ROUGE often exhibit weak correlation with human preferences~\citep{stiennon2022learning}, and previous studies have shown that fine-tuning language models with PPO using human preferences leads to more effective summaries. We compare different methods by generating completions on the test split of the TL;DR summarization dataset, then calculating the average win rate against reference completions from the test set. Completions for all methods are generated across a range of temperatures from 0.0 to 1.0, with win rates presented in Figure~\ref{fig:frontier-tldr-main} (right). DPO, PPO, and Preferred-FT all fine-tune the identical GPT-J SFT model\footnote{\url{https://huggingface.co/CarperAI/openai_summarize_tldr_sft}}. We observe that DPO achieves a win rate around 61\% at a temperature of 0.0, outperforming PPO, which attains approximately 57\% at its optimal temperature of 0.0. Additionally, DPO reaches a higher peak win rate compared to the best baseline among the top $N$. It is important to mention that we did not extensively tune DPO's $\beta$ hyperparameter, so the reported performance might underestimate its full capability. Furthermore, DPO demonstrates greater robustness to changes in sampling temperature than PPO, whose performance can decline to that of the base GPT-J model at elevated temperatures. Preferred-FT shows little improvement over the SFT model. A direct comparison between DPO and PPO in human evaluations is also presented in Section~\ref{sec:human-judgments}, where DPO samples generated at temperature 0.25 were favored 58\% of the time over PPO samples generated at temperature 0.

For single-turn dialogue, we assess various approaches using a subset of the test portion of the Anthropic HH dataset \citep{bai2022training} involving one step of human-assistant interaction. GPT-4 evaluations utilize the preferred completions from the test set as references to calculate the win rate across different methods. Since there is no established SFT model for this task, we begin with a pre-trained Pythia-2.8B, apply Preferred-FT to train a reference model on the selected completions so that these completions fall within the model's distribution, and subsequently train with DPO. We also compare against the top result from 128 Preferred-FT completions (noting that the Best of $N$ baseline plateaus at 128 completions for this task; see Appendix Figure~\ref{fig:best-of-n}) and a 2-shot prompted variant of the Pythia-2.8B base model. Our findings indicate that DPO matches or surpasses the best-performing temperatures of these methods. Additionally, we evaluate an RLHF model trained with PPO on the Anthropic HH dataset\footnote{\url{https://huggingface.co/reciprocate/ppo_hh_pythia-6B}} from a reputable source\footnote{\url{https://github.com/CarperAI/trlx/tree/main/examples/hh}}, but fail to identify any prompt or sampling temperature that outperforms the base Pythia-2.8B model. Given our TL;DR results and the fact that both approaches optimize the identical reward function, we treat Best of 128 as an approximate proxy for PPO-level performance. Overall, DPO stands out as the only computationally efficient method that improves upon the preferred completions within the Anthropic HH dataset, delivering comparable or superior results relative to the computationally intensive Best of 128 baseline. Lastly, Figure~\ref{fig:dialogue-main} illustrates that DPO attains its peak performance relatively quickly.

\subsection{Generalization to a novel input distribution}

To further assess the comparative performance of PPO and DPO under shifts in data distribution, we test the PPO and DPO policies obtained from our Reddit TL;DR summarization experiment on a different dataset: news articles from the test portion of the CNN/DailyMail dataset \citep{nallapati-etal-2016-abstractive}, employing the optimal sampling temperatures identified on TL;DR (0 and 0.25). The outcomes are shown in Table~\ref{tab:ood}. We calculated the GPT-4 win rate against the original summaries in these datasets, utilizing the same GPT-4 (C) prompt used for Reddit TL;DR, with the phrase ``forum post'' replaced by ``news article''. For this alternative distribution, DPO continues to outperform PPO by a large margin. This experiment offers preliminary evidence that DPO policies can generalize at least as effectively as PPO policies, despite DPO not leveraging the extra unlabeled Reddit TL;DR prompts that PPO utilizes.

\subsection{Validating GPT-4 assessments through human judgments}
\label{sec:human-judgments}

We perform a human evaluation study to confirm the reliability of GPT-4's assessments, drawing on the results from the TL;DR summarization experiment and two distinct GPT-4 prompts. The \textbf{GPT-4 (S)} (simple) prompt asks solely which summary better captures the important information in the post. The \textbf{GPT-4 (C)} (concise) prompt additionally requests which summary is more concise; this prompt is evaluated because we observe that GPT-4 tends to prefer longer, more repetitive summaries than humans do when using the \textbf{GPT-4 (S)} prompt. Complete prompt details are provided in Appendix~\ref{app:prompts}. We conduct three comparisons involving the highest-performing method (DPO, temp. 0.25), the lowest-performing method (PPO, temp. 1.0), and an intermediate method (SFT, temp. 0.25) to represent a range of sample qualities; all three methods are compared against greedily-sampled PPO (its top-performing temperature). Results indicate that GPT-4, with both prompts, tends to concur with human judgments roughly as often as humans agree amongst themselves, suggesting GPT-4 is a reasonable surrogate for human evaluation (due to a limited number of human raters, multiple human judgments were gathered only for the DPO versus PPO-1 comparisons). Overall, the \textbf{GPT-4 (C)} prompt yields win rates that more closely align with human assessments; hence, we employ this prompt for the primary results in Section~\ref{sec:dpo-real-datasets}. Additional information about the human study, including the rater interface and list of volunteers, can be found in Appendix~\ref{app:human-study}.

\section{Discussion}
Learning from preferences represents a robust and scalable approach for developing capable and aligned language models. We have presented DPO, a straightforward training method that enables language models to be trained from preferences without relying on reinforcement learning. Instead of forcing the preference learning task into a conventional RL framework to leverage existing RL algorithms, DPO establishes a direct correspondence between language model policies and reward functions. This connection allows a language model to be trained to meet human preferences \textit{directly} using a simple cross-entropy loss, eliminating the need for reinforcement learning and avoiding any loss of generality. With minimal hyperparameter tuning, DPO achieves performance on par with or surpassing current RLHF methods, including those based on PPO; consequently, DPO significantly lowers the barrier for training language models from human preferences.

\textbf{Limitations \& Future Work.} Our findings prompt several important avenues for future investigation. How well does the DPO policy generalize to out-of-distribution data compared to approaches that optimize an explicit reward function? Preliminary evidence indicates that DPO policies generalize comparably to PPO-trained models, though a more thorough analysis is warranted. For instance, is it possible for training with self-labeling derived from the DPO policy to effectively utilize unlabeled prompts? Another consideration is how reward over-optimization appears within the direct preference optimization framework, and whether the slight performance decline shown in Figure~\ref{fig:dialogue-main}-right exemplifies this phenomenon. Furthermore, while our experiments cover models up to 6B parameters, investigating the scalability of DPO to state-of-the-art models that are orders of magnitude larger represents a promising research direction. Regarding evaluation methodologies, we observe that GPT-4's computed win rates are influenced by the prompt used; future work could explore optimal strategies to elicit reliable and high-quality judgments from automated evaluators. Lastly, the potential applications of DPO extend beyond training language models based on human preferences, including the training of generative models in other data modalities.

\end{document}